\documentclass[11pt]{article}
\usepackage[margin=1in]{geometry}
\usepackage{listings}
\usepackage{color}
\usepackage{graphicx} % Required for inserting images

\title{COSC 417 Programming Task 1 \\ DFA Simulator}
\author{Jonathan Llovet, Matt Saxton, Amarachi Ufelle King, Abdullah Radid}
\date{2nd April 2026}

\definecolor{codegray}{gray}{0.9}

\lstset{
    backgroundcolor=\color{codegray},
    basicstyle=\ttfamily\small,
    breaklines=true,
    frame=single
}

\begin{document}

\maketitle

\section*{Program Description}

This project aims to simulate a Deterministic Finite Automaton (DFA) using the Java programming language. The program reads a file containing the machine's description and builds the machine as a function on a given input string; it then decides whether the string is to be accepted or rejected.

\subsection*{Overview}

The machine is defined as a DFA, $M$, of the form, $M = (Q, \Sigma, \delta, q_0, F)$, where:
\begin{itemize}
    \item $Q$ represents a finite set of states, each representing the effect of an input symbol on the configuration of the machine:
    $1 \leq |Q| \leq 20$, where $20$ is implemented as a constant integer \texttt{MAX\_STATES}. States are numbered $1$ through $n$ as specified by the machine description.
    \item $\Sigma$ represents a finite set of symbols or alphabet that defines what inputs are valid and can be read by the machine: $\Sigma = \{0, 1\}$, enforced by an If-statement and \texttt{IllegalArgumentException}. While this DFA is hard-coded to only recognize binary input ($|\Sigma| = 2$, implemented as a constant integer \texttt{ALPHABET\_SIZE}), it can convert any 2-symbol alphabet to $\Sigma$.
    \item $\delta$ represents the transition function, the means by which the machine changes from one state to another upon reading a valid input symbol; for each state and input symbol, there is one and only one next state: $Q \times \Sigma \rightarrow Q$, implemented as a 2-dimensional integer array \texttt{transition}, having parameters [\texttt{MAX\_STATES}] and [\texttt{ALPHABET\_SIZE}].
    \item $q_0$ represents the starting state, the initial configuration of the machine: $q_0 = \{1\}$.
    \item $F$ represents the accepting states; if the machine is found to be configured in one of these states once it has finished reading a string of input symbols, the string is considered accepted by the DFA: $F \subset Q$ as specified by the machine description, implemented as a \texttt{Hashset$<$Integer$>$} to prevent duplicate accepting state declarations and for fast (O(1)) accepting/non-accepting final state determination.
\end{itemize}

\newpage
\subsection*{File Input Format}

The DFA is read from a text file consisting of only integer values with the following structure:
\begin{enumerate}
    \item First line: number of states $n$
    \item Second line: accepting states
    \item Remaining lines: transitions of the form
    \[
    (\texttt{current state, input symbol, next state})
    \]
\end{enumerate}
Note that numbers on the same line are separated by spaces.
Example:
\begin{verbatim}
5
2 4
1 0 2
1 1 4
...
\end{verbatim}

\subsection*{Data Structures Used}

\begin{itemize}
    \item \textbf{Integer (int)}: Stores the number of states in range $1 < n <= 20$. Alphabet validation is performed with raw integers 0 and 1.
    \item \textbf{Hashset$<$Integer$>$}: Stores accepting states for fast lookup using \texttt{contain()}.
    \item \textbf{2D Array (int[][])}: Implements the transition function.
    Specifically:
    \[
    \texttt{transition[state][symbol]}
    \]
    where:
    \begin{itemize}
        \item \texttt{state} is the current state (1 to 20)
        \item \texttt{symbol} is either 0 or 1
    \end{itemize}
    \item \textbf{String}: Represents the input string, whose elements are read as integers to compare with binary alphabet.
\end{itemize}

\subsection*{Simulation Runtime Logic}

\begin{itemize}
    \item The simulation starts at state 1.
    \item Each symbol in the input string is processed sequentially.
    \item The next state is determined using the transition table.
    \item If an undefined transition is encountered (state 0), the string is rejected.
    \item After processing the input:
    \begin{itemize}
        \item If the final state is in $F$, output \texttt{ACCEPT}
        \item Otherwise, output \texttt{REJECT}
    \end{itemize}
\end{itemize}

\newpage

\section*{Java Code}

\begin{lstlisting}[language=Java]
% copy java code here
\end{lstlisting}

\newpage

\section*{DFA}

\subsection*{File Contents}
DFA 1 (Example 1.11, page 38):
\begin{verbatim}
5
2 4
1 0 2
1 1 4
2 0 2
2 1 3
3 0 2
3 1 3
4 0 5
4 1 4
5 0 5
5 1 4
\end{verbatim}

\noindent DFA 2 (Example 1.68(a), page 76):
\begin{verbatim}
3
2 3
1 0 2
1 1 3
2 0 1
2 1 2
3 0 2
3 1 1
\end{verbatim}

\subsection*{Test Cases}
\begin{itemize}
    \item Accepted:
        010 (DFA 1), % rule is "start and end with same symbol"
        001 (DFA 2) % rule is "must have an odd number of 0s, or end in an odd number of 1s" - note 0=a, 1=b
    \item Rejected:
        011 (DFA 1), % rule is "start and end with same symbol"
        100 (DFA 2) % rule is "must have an odd number of 0s, or end in an odd number of 1s"
\end{itemize}


\newpage

\section*{Program Output Screenshots}

\begin{figure}[h!]
    \centering
    \includegraphics[width=0.8\textwidth]{screenshot1.png} % put image file name here
    \caption{Sample Run - DFA 1}
\end{figure}

\begin{figure}[h!]
    \centering
    \includegraphics[width=0.8\textwidth]{screenshot2.png} % put image file name here
    \caption{Sample Run - DFA 2}
\end{figure}


\end{document}
